\section{General Discussion}
\label{GenDiscussion}
\setlength{\parindent}{20pt}
\justifying

\subsection{\large{\emph{Creation of a novel implicit bilingual SL task}}}
\noindent
The present project followed the development of an implicit bilingual statistical learning task, in which the artificial words presented to participants were formed of morphemes that co-occurred only with other morphemes of that language. The task for participants was to segment these words to learn the morphemes, and track these co-occurrence patterns to be able to accurately reject between-language combinations presented in the testing phase. This two-pronged task therefore called for both extraction and integration using the frequency-based and co-occurrence statistical information present in the artificial stimuli.

% morph fam effect
This paradigm allowed for the replication of the morpheme familiarity effect found in \textcite{Lelonkiewicz2020}. Participants were able to develop a sensitivity to morphemes in novel test words following only a passive word exposure phase. This uptick in their feeling of congruency going from test strings with no known morphemes to those with one morpheme present, to strings composed of two known morphemes was a very solid phenomenon observed across all three experiments. This was the case even for the first experiment, in which the design made the task the most difficult. 

Interestingly, the experiment 3 stimuli font change to Latin characters had the consequence of reducing the proportion of ``yes'' responses across all conditions. The one-morpheme items received affirmative evaluations half of the time, which nicely reflects their two possible interpretations. The presence of a foreign element within the test string may cause for its rejection, or the recognition of a familiar element might be enough for the participant to generalise the learned patterns and accept a novel pairing, much as an English learner knowing ``read$\vert$er'' might accept ``read$\vert$ing''.

The morpheme familiarity effect was reinforced based on multilingual factors. Participants who made more use of their second rather than their first language had greater separation of their responses to the different testing  conditions (in experiments 2 and 3). The balanced use of an L1 and L2, or the greater use of an L2 over an L1 would be more cognitively effortful to allow for communication in the weaker language. However, over time, this may have led to greater linguistic metacognitive processing as they become more bilingual (\cite{Ordin2020}). This was also the case for those with more bilingual balance (more similar scores for their L1 and L2; experiments 2 and 3). As it has been proposed that bilingual balance has links to language processing abilities (\cite{Birdsong2016}), these participants may be more efficient at morphemic identification. This evidence would support the theory that those with more bilingual (as opposed to monolingual) profiles have built up greater sensitivity to familiar chunks. Additionally, in experiments 1 and 3, trilinguals were the group with the best morpheme familiarity, followed by bilinguals, with quadrilinguals' performance falling below even that of monolinguals. This would indicate that there is at least a bilingual morphemic sensitivity effect that may be reversed beyond 3 languages.

% sig group effect
Within the two-morpheme condition, the experimental conditions in the final version were sufficient to find a significant effect of implicit bilingual co-occurrence learning even at a group level. Participants were therefore able to track co-occurrences across time and the different training blocks in which the morphemes were presented to infer that this web of pairings denotes items coming from different languages. These statistical regularities learned completely unconsciously were sufficient to create mental categorisation of the input. 

% complex task: extraction then integration
While the size of the effect is very modest, this is an ability that was significantly above chance level. This design is both novel and complex, which may have contributed to the difficulty in finding more decisive results. Indeed, participants are required to first engage in extractive statistical learning to find the morphemes within the words, and successively integrate the morphemes into categories thanks to co-occurrence tracking mechanisms. There may be ways to facilitate the segmentation, but the integration seems to be the trickier process to manipulate. None of the multilingual metrics reliably predicted familiarity scores across multiple experimental versions, and none did so in experiment 3. This may have been anticipated, as \textcite{Poepsel2016} suggested that a bilingual advantage may only exist for integrative statistical learning.

The question remains as to how to make the integration process easier, while keeping the statistical learning manipulation implicit. More cognitive resources might simply be freed up for integration thanks to conditions supporting easier segmentation, but the unknowns behind the process of co-occurrence itself limit the extent to which this can be directly observed. The statistical regularities are already fully non-overlapping across languages (that is, a morpheme from language 1 was never shown with one from language 2 during training), so this categorisation cue was fully reliable. Other experimental manipulations to make this cue more salient unfortunately sacrifice the desired implicitness of the training.

% interpretation of SL structure
However, the interpretation of this implicit co-occurrence structure by participants may not have been the one conceived of by experimenters during the design process. Despite the lack of biasing instructions favouring the interpretation that the input was monolingual, all of the testing words may have been learned as if they do come from a single language. Consider the below table of mock stimuli in English:

\begin{table}[!htbp]
    \begin{center}
        \begin{tabular}{@{}lll@{}}
            \toprule
            & \textbf{``L1''} & \textbf{``L2''} \\ \midrule
            congruent & sting - er & walk - ing \\ \midrule
            congruent & weak - est & quick - ly \\ \midrule
            & & \\ \midrule
            ``incongruent'' & weak - ly & walk - er \\ \midrule
            ``incongruent'' & sting - ing & quick - est \\ \midrule
        \end{tabular}
    \end{center}
\end{table}

In this ``L1'', two stems are associated with ``-er'', a suffix indicating a noun part of speech, and ``-est'', which indicates an adjective or adverb. ``L2'' has the suffixes ``-ing'', for verbs, and ``-ly'', for adverbs. However, from looking at the ``incongruent'' cross-language combinations, it is perfectly plausible for one stem to be associated with different suffixes, and even unusual combinations of known stems and suffixes would not necessarily be cause for rejection (``giraffing'' could be a form of animal play-acting similar to monkeying). Therefore, instead of treating a ``between-language'' combination as a violation, it could be understood by a similar explanation as a form of generalisation of learned morpheme part of speech indicators. However, even if this was the case in the present experiment, this would still require the same mental integration of classifying linguistic stimuli as different parts of speech, and would still involve co-occurrence tracking.

\subsection{\large{\emph{Multilingual variables}}}
\noindent
% strength of experiment
One of the strengths of the present experiment was the ability to recruit participants from a very wide variety of language backgrounds. Through the use of artificial stimuli, and especially with the BACS-1 characters from the first two experiments, the authors were able to include participants from any geographical region and speaking any languages, with the only requirement being English fluency to understand task instructions. The findings outlined here are therefore not restricted to specific stimulus language combinations or to specific bilingual participant profiles. This allows for more generalisability of the statistical learning effects found. Wide-ranging artificial statistical learning studies may be a way to overcome the anglocentrism and eurocentrism that characterise most studies today.

% sig vars
Several conclusions can therefore be drawn about the effect of multilingual variables in general with relation to statistical learning in this task. The only variable that predicted co-occurrence tracking ability was multilingual balance, as calculated by cosine similarity. The more balanced the multilingual profile across up to four spoken languages, the worse they seemed to be at forming these statistical regularity-based categories. Past research had found a role for increased \emph{bilingual} balance providing a benefit in a different bilingual statistical learning paradigm (\cite{Onnis2018}). However, there was no effect of bilingual balance in this experiment, meaning that the cause here must have multilingual roots in effects beyond two languages. Additionally, there was no effect on multilingual experience on co-occurrence tracking: the mere accumulation of linguistic knowledge did not play a role. Looking across all versions of the experiment, the bi-/tri-/quadrilingual categories did not have a consistent effect either; the cosine similarity significance therefore cannot be attributed to multilingual status, but seems to have to do with how third and fourth languages are being learned, used, and valued.

The participants' age range was set as 18-25 in order to remove any age-related statistical learning differences. However, being perfectly proficient or balanced in four languages at that young is perhaps a difficult thing to be: having equal capacities in four languages at that age might mean shallower learning of all (or most) of them. This phenomenon may lessen for older balanced multilinguals who have had more time to solidly learn all of their languages. An earlier acquisition, exaggerated use, or excessive valuing of a third and fourth language might have impacts on the ability to perform integrative statistical learning to form distinct categories. \textcite{Gullifer2019} posited that bilinguals are able to perform in a compartmentalised or integrated way. Additionally, \textcite{Abutalebi2016} proposed that, according to the adaptive control hypothesis, the individual differences that change show bilinguals represent, access, and control their languages might determine how integrated their use is. It would stand to reason that the same can be extended to multilinguals. Those multilinguals who are more used to having their other languages co-activated in various settings may more easily integrate statistical regularities across languages.

% cossim
As for the measure of cosine similarity more generally, it was initially developed to overcome the perceived mathematical correlation between entropy and multilingual experience. Empirically, a very high correlation was seen across all versions of the experiment: experiment 1 Pearson correlation t(191) = 37.30, p $<$ 2.2e$^{-16}$, estimate = 0.94; experiment 2 Pearson correlation t(182) = 30.18, p $<$ 2.2e$^{-16}$, estimate = 0.91; experiment 3 t(193) = 32.95, p $<$ 2.2e$^{-16}$, estimate = 0.92. However, the correlation between cosine similarity scores and multilingual experience are only tentative: experiment 1 Pearson correlation t(191) = -1.68, p = 0.09, estimate = -0.12; experiment 2 Pearson correlation t(182) = -1.62, p = 0.11, estimate = -0.12; experiment 3 Pearson correlation t(193) = -1.46, p = 0.14, estimate = -0.10. The correlation between entropy and cosine similarity also demonstrate this: experiment 1 Pearson correlation t(191) = -3.60, p = 0.0004, estimate = -0.25; experiment 2 Pearson correlation t(182) = -4.85, p = 2.63e$^{-6}$, estimate = -0.34; experiment 3 Pearson correlation t(193) = -5.11, p = 7.87e$^{-7}$, estimate = -0.34. Entropy therefore seems to be a measure at the crossroads between balance and multilingual experience, though being very largely biased by the latter. Cosine similarity is the more independent measure of the two.

However, some issues still remain surrounding the practicalities of its use. Given the calculation of balance either with global language scores or by conserving only proficiency or use, it is sometimes difficult to see patterns across literature. The calculation of cosine similarity can be undergone in the same way regardless of the choice made, but consensus on this point is greatly needed. Another point is that of possible cut-offs put in place in deciding which global language scores to accept in the calculation of cosine similarity. Participants were instructed to enter any number of languages they speak, but the choice to include or omit some very minimal understanding of one language may have important consequences on the calculated cosine similarity score. If a bilingual has the maximal score in their L1 and L2 (score vector = [218, 218]), they will get a cosine similarity of 1, representing perfect balance. However, if they elected to also mention their very dim knowledge of an L3 (score vector = [218, 218, 10]), their cosine similarity would fall to 0.83. On a scale from 0 to 1, this is a very significant difference. For this reason, the present authors chose as 50 the minimum global language score: this represents 23\% of the perfect global language score, and we esteemed that this showed sufficient additional learning and knowledge to add to a multilingual's linguistic profile. However, this might be too strict, or too liberal; it is a first attempt to create a reasonable, representative measure of multilingual balance. The one other issue with this measure is the small variability: due to the nature of the global language scores, this number is necessarily positive, and restricted even more due to the cut-off. Across all experiments, the lowest cosine similarity score was 0.84. Another question can be asked of the monolinguals. By default, since they only speak one language, they have a cosine similarity of 1, but it cannot be said that this one-dimensional calculation represents balance. It would stand to reason to exclude monolinguals from cosine similarity analyses to preserve the true meaning of the measure.

% PCA
In terms of what other multilingual measures might be particularly helpful in analyses, the Principal Components Analyses done often picked up similar components across all experiments. There were always L3 and L4 components, as well as L1 - L2 use and L2 history. Their prevalence might suggest that these components are good bottom-up summaries of the multilingual experience. While L3 and L4 components have their natural place in research, L2 history (which charts answers to questions such as L2 age of acquisition) might explain more variance than previously thought. Theoretically, the length of time that one has juggled more than one language might be very significant in understanding the multilingual experience. The split between L1 and L2 use is also very theoretically compelling: the extent to which one stays in their dominant language or expends more cognitive effort while developing capabilities in a second language might be a component to preserve in future research. 

% category
Finally, the overall multilingual category that participants belong to does seem to be a useful catch-all to understand global effects. However, any general effect of being a bilingual or trilingual remains uninformative as to which aspect of bi- or trilingualism provoked said effects. It may be the starting point for further exploration into multilingual use, balance, or experience. However, more sensitive measures are definitely required to look specifically at language history, use, proficiency, attitudes, and balance in all of a participant's languages to understand the complex dynamics of multilingualism.


\subsection{\large{\emph{Limitations}}}
\noindent
In terms of the experimental design, many elements may require tweaking to render the experiment more easily digestible to the statistical learner. Passive exposure during the training phase may not be engaging enough to sustain attention throughout the length of the training. This might also not be realistic, since real children interact with their caretakers and get feedback during learning (\cite{Frost2019}). There may also not be enough repetitions of the stimuli during training to give the participant sufficient time to segment and track co-occurrences, especially in experiments 1 and 3, which did not have the morpheme spelling task to aid their segmentation. The number of individual morphemes used in the experiment might also be too high, which would complicate tracking. The modifications of experiment 3 did seem to make learning smoother for participants, but it is regrettable that this limited the participant backgrounds due to the necessity to validate the experimental stimuli. 

The usage of 2-artificial forced-choice testing also only permits post-hoc statistical learning checks, not online learning measures to examine the timeline. It might have been interesting to see after which number of repetitions the morphemes were learned well enough to divert resources to co-occurrence tracking. Evidence of this learning might also be seen earlier in neurological measures before they emerge in behaviour (\cite{Aguasvivas2024}). It might also be possible through electrophysiological measures to detect larger differences in  between- and within-language two-morpheme combinations. The statistics presented in statistical learning experiments are also very watered-down versions of the many regularities learners are bombarded with in their ordinary lives, and are rarely as neat (\cite{Frost2019}). Real bilingual input would indeed also have both implicit and explicit cues to differentiate input from different languages (\cite{Weiss2009}).


\subsection{\large{\emph{Future directions}}}
\noindent
This statistical learning task is still in its infancy, and more marked effects might be seen thanks to some refinements. Further piloting could help identify the best parameters for the number of repetitions of the initial words in training, the number of distinct morphemes used, and so on. The 1M and 2M conditions might be removed to reduce the strong ``yes'' bias in the 2M condition and allow for participants to reply on finer instincts than the mere presence of two familiar morphemes. Additionally, due to the presence of BACS-1 characters in experiments 1 and 2, participants had no phonological information available to support their learning (\cite{Chetail2017}). This was, to a lessen degree, also the case in experiment 3, where approximately half of the morphemes were just composed of consonants, and had low pronounceability (e.g. ``VWNM''). These items were included to keep a large pool of morphemes for participants to be randomly assigned; in this way, having different morphemes for each iteration of the experiment would greatly reduce between-participant stimuli-related confounding factors. However, the morphemes with a vowel may be enough to satisfy this condition in a future experiment, and boost learning through the addition of phonology.